% !TeX root = ../../Book3.tex
\chapter{Riemann 映照定理}
\label{b3:ch:05}
\BookChapterNumber{b3:ch:05}
\begin{chapterabstract}
Riemann 映照定理说明，每个不等于全平面的单连通平面区域都有双全纯的圆盘坐标。存在性的证明先利用全纯平方根构造一个有界单叶函数，再在固定基点和导数相位的函数族中最大化导数。Montel 紧性与单叶极限定理使极值可达，遗漏点的平方根变换则迫使极值映射满射。Schwarz 引理给出归一化唯一性，也把逆映射的导数解释为区域在基点处的共形半径。显式映射、区域穷尽及全平面的例外共同说明这个结论的适用范围。
\end{chapterabstract}

\begin{learninggoals}
\begin{itemize}
\item\label{b3:item:05-goal-1}能从原函数或闭路周期判断全纯对数与根的存在，并区分某个函数有根与整个区域单连通这两种性质。
\item\label{b3:item:05-goal-2}能完整说明归一化单叶族非空、极值可达和极值映射满射的证明，指出真区域与非退化条件各自的用途。
\item\label{b3:item:05-goal-3}能计算遗漏值平方根变换的导数放大因子，并在显式映射中同时核对分支、整体像和导数相位。
\item\label{b3:item:05-goal-4}能用 Schwarz 等号情形证明归一化唯一性，计算共形半径，并理解区域扩大或穷尽时归一化映射的变化。
\end{itemize}
\end{learninggoals}

半平面有显式的分式线性圆盘坐标，角域可以先取一个幂函数，水平带可以先作指数变换。这些例子通常从边界形状猜出函数，再逐一验证它在内部单射且满射。面对不规则边界，这种猜公式的方法可能无从开始。Riemann 映照定理回答一个更一般的存在问题：只要区域单连通且不是全平面，无论能否写出初等公式，都能建立一张双全纯的圆盘坐标。

\ProseParagraphBreak

这里的结论要求全局可逆。导数处处非零只保证局部可逆，例如指数函数在全平面有周期；把边界的若干点送到圆周，也没有证明整个内部恰好映到圆盘。反过来，定理本身只建立区域内部的双全纯等价，并不承诺映射能够连续延到欧氏边界。分支的单值性、内部的单射与满射、指定基点处的归一化，以及边界延拓，承担着不同的数学要求。

\ProseParagraphBreak

本章所需的紧性已经在第四章建立。\cref{b3:thm:04-montel}用于从取值于圆盘的函数列抽取局部一致子列；\cref{b3:thm:04-univalent-limit}说明单叶极限若不为常数就仍单叶，而一个正的极限导数将排除常数退化。另一方面，\cref{b3:thm:02-schwarz}对固定原点的圆盘自映射给出导数模不超过1，并在取等时确定它是旋转。这些结果必须连同各自条件一起使用，才能从一个全纯极限得到满足归一化的双全纯映射。

\ProseParagraphBreak

先在\secref{b3:sec:0501}准备对数与根，再在\secref{b3:sec:0502}完成极值构造，最后在\secref{b3:sec:0503}处理唯一性、共形半径和具体映射。存在证明中的两次取平方根有不同用途：第一次制造一个被像集避开的小圆盘，以便取有界倒数；第二次针对假设存在的遗漏值，提高已归一化的导数。Theodore W. Gamelin 的《Complex Analysis》\cite{Gamelin2001Complex}和 Lars Ahlfors 的《Complex Analysis: An Introduction to the Theory of Analytic Functions of One Complex Variable》\cite{Ahlfors2021Complex}可供继续阅读共形映射理论。

% !TeX root = ../../Book3.tex
\section{单连通区域}
\label{b3:sec:0501}

把一个区域映成圆盘，是寻找全局单射函数的问题。局部逆函数定理只保证小邻域可逆，不能给出整域的共同坐标。单连通性使无零函数可取对数与平方根；这些分支构造将产生足够多的有界单叶函数。

\subsection{单连通性}
\label{b3:subsec:050101}

本章单连通性沿用\subsecref{b3:subsec:020203}的闭路径收缩定义，区域必须非空、开且连通。凸区域、星形区域单连通；单位圆盘去掉一条从边界伸入内部的径向线段仍单连通，去掉一个内部点则不再单连通。

\ProseParagraphBreak

平面区域的拓扑判断不等于观察一幅图。一个常用的平面拓扑判据是：区域在扩充平面中的补集连通，当且仅当它单连通。该判据涉及 Jordan 分离与平面路径的多边形逼近，可作为本章所需的平面拓扑基础；实际分析证明仍以 Cauchy 定理的收缩条件为准。本章的构造不假定边界光滑，也不要求区域有界或边界只有一条简单曲线。

\ProseParagraphBreak

全平面本身单连通，却不能双全纯映到圆盘。否则圆盘坐标作为有界整函数由 Liouville 为常数，违反单射性。因此存在定理必须排除 \(U=\mathbb C\)，这一排除将具体用于在区域之外选择一点。

区域不必有简单边界。例如平面减去非负实轴是单连通的：极坐标可取 \(0<\arg z<2\pi\)，从而用 \(\log z\) 把它映到一条水平带，带是凸的。平面减去两个点却有两个独立绕行，绕其中任意一个点的闭路都不能收缩。对带有许多边界分支的区域，直接寻找一个能够连续选择的坐标，比只数图上的“孔”可靠。边界有尖点也不妨碍单连通性；是否能把共形映射连续延到尖点，属于第7章的另一问题。

\ProseParagraphBreak
单连通性只在区域内部讨论。开圆盘和闭圆盘不是同一对象，后者不是平面区域；开圆盘去掉一个边界点仍是原来的开圆盘。这个区分在谈到从边界伸入的裂缝时尤为必要。删去的线段应当确实含有内部点，其边界端只负责把障碍连接到区域外部，不能作为内部点参加闭路径。

\subsection{全纯对数}
\label{b3:subsec:050102}

\begin{theorem}[全纯对数]\label{b3:thm:05-holomorphic-log}
设 \(U\) 为单连通区域，\(f\) 在 \(U\) 上全纯且处处非零。则存在 \(U\) 上全纯函数 \(L\)，使 \(e^L=f\)。任意两个这样的函数相差 \(2\pi\mathrm{i}k\)，其中 \(k\in\mathbb Z\)。
\end{theorem}
\begin{proof}
函数 \(f'/f\) 全纯，由单连通 Cauchy 定理有原函数 \(H\)。求导得 \((f e^{-H})'=e^{-H}(f'-fH')=0\)，所以 \(fe^{-H}=c\ne0\)。选一个复数 \(\ell\) 满足 \(e^\ell=c\)，令 \(L=H+\ell\)，便有 \(e^L=f\)。若 \(L_1,L_2\) 都满足条件，差的指数为1，故差值落在离散集合 \(2\pi\mathrm{i}\mathbb Z\)。连续性与连通性使差为一个固定整数倍。
\end{proof}
给定 \(a\in U\) 和一个允许的 \(L(a)\)，构造中的常数选择就唯一。这个归一化对于沿路径继续分支很有用。

\ProseParagraphBreak

无零条件必不可少，因为指数函数从不为零。单连通性也不能直接删除：\(f(z)=z\) 在去心平面无零，却无全局对数。但它不是存在对数的必要条件，例如常值1在任何区域都有对数0；对某个特定函数的精确条件是 \(f'/f\) 的所有闭曲线周期为零。

\subsection{全纯平方根}
\label{b3:subsec:050103}

\begin{corollary}[全纯根]\label{b3:cor:05-holomorphic-root}
设 \(f\) 在单连通区域 \(U\) 上全纯且无零，\(m\ge1\) 为整数。则 \(f\) 有全纯 \(m\) 次根；任意两个根相差一个固定 \(m\) 次单位根因子。
\end{corollary}
\begin{proof}
取前一定理的 \(L\)，则 \(g=e^{L/m}\) 满足 \(g^m=f\)。两个无零根之商取值于有限集合 \(\{\zeta:\zeta^m=1\}\)，连续性与连通性使其恒定。
\end{proof}
若允许零点，局部 \(m\) 次根要求每个零点阶数能被 \(m\) 整除；例如 \(z\) 在含0的圆盘没有全纯平方根，\(z^2\) 却有。全局还需核对去掉零点后沿闭路的周期。下节只对无零函数取平方根，所以无需把这一更一般的根判别当作前提。

\ProseParagraphBreak

对 \(z-b\) 的平方根，若 \(b\notin U\)，得到 \(h^2=z-b\)。由 \(h(z_1)=h(z_2)\) 得 \(z_1=z_2\)，所以根本身单射。它的像与其负像不相交：若 \(h(z_1)=-h(z_2)\)，平方后点相同，再得 \(h(z_1)=0\)，与 \(b\notin U\) 矛盾。这个像集分离性质将产生有界函数。

\begin{proposition}\label{b3:prop:05-roots-with-zeros}
设 \(U\) 单连通，\(f\) 在 \(U\) 全纯且不恒零，\(m\ge2\)。则 \(f\) 有全纯 \(m\) 次根，当且仅当每个零点的阶数都能被 \(m\) 整除。
\end{proposition}
\begin{proof}
必要性由局部展开 \(f=g^m\) 立即得到。证明充分性时先去掉离散零点集 \(Z\)。对 \(U\setminus Z\) 中任意闭曲线，单连通 Cauchy 定理和有限奇点的留数公式给
\[
 \frac1{2\pi\mathrm{i}}\int_\gamma\frac{f'}f\,\mathrm{d}z
 =\sum_{a\in Z}\operatorname{Ind}(\gamma,a)\operatorname{ord}_a(f).
\]
右侧只有有限项且为 \(m\) 的整数倍。这里对于非简单曲线，可先按有限多边形分割成简单圈；小扰动避开零点后局部原函数保证积分不变。因此沿路径定义 \(g=c\exp(\frac1m\int f'/f)\)，其中 \(c^m=f(z_0)\)，闭路积分的指数均为1，得到单值全纯根。在零点 \(a\) 附近写 \(f=(z-a)^{mk}u\)，\(u\) 无零且在小盘有根 \(v\)。穿孔小盘上 \(g/((z-a)^k v)\) 是固定单位根，因此 \(g\) 全纯延过 \(a\)。
\end{proof}
与无零函数相比，证明多了两个步骤：先在删去零点的区域检查周期，再在每个零点局部补回。后一个步骤不能替代前一个步骤。在非单连通区域，所有零点阶数合适也可能没有根，例如去心平面的 \(z\) 没有零点，却无平方根。

\subsection{单叶函数}
\label{b3:subsec:050104}

区域上的单射全纯函数称为\term[danye-hanshu]{单叶函数}[univalent function]。单叶函数的导数不为零，由第一章局部标准形可知；其逆在像区域上全纯，因此是双全纯映射。本卷的“共形映射”若表示两个区域之间的等价，明确指这种双全纯映射，而局部共形函数不必全局单叶。

\ProseParagraphBreak

例如 \(e^z\) 导数处处非零，在宽度小于 \(2\pi\) 的水平带内单叶，却在整个平面不单叶。函数 \(z^2\) 在右半平面单叶，在包含0的圆盘不单叶，而且在0导数为零。判断单叶性需要比较任意两点，不能只核对导数。

\ProseParagraphBreak

若单叶函数 \(f:U\rightarrow\mathbb D\) 在某点 \(a\) 已确定值，可用 \(\phi_{f(a)}(w)=(w-f(a))/(1-\overline{f(a)}w)\) 将其送到0，再乘单位复数使导数成为正实数。这种归一化消除圆盘自同构中的自由度，不改变像是否为整个圆盘。

% !TeX root = ../../Book3.tex
\section{极值方法}
\label{b3:sec:0502}

存在性的证明不直接解一个映射公式，而是把原点附近的伸缩率推到最大。正规族保证极值能达到；若极值映射漏掉一点，就用平方根重新铺开像，使伸缩率变大。两项共同给出满射性。

\subsection{归一化族}
\label{b3:subsec:050201}

固定真单连通区域 \(U\subsetneq\mathbb C\) 与 \(a\in U\)。令 \(\mathcal F\) 为所有单叶全纯映射 \(f:U\rightarrow\mathbb D\)，满足 \(f(a)=0\)、\(f'(a)>0\)。
\begin{lemma}\label{b3:lem:05-nonempty-normalized}
归一化族 \(\mathcal F\) 非空。
\end{lemma}
\begin{proof}
选 \(b\notin U\)，取 \(h(z)^2=z-b\) 的全纯根。它单叶，像与负像不交。由于像在 \(h(a)\) 附近包含一个小圆盘 \(D(h(a),r)\)，像便避开 \(D(-h(a),r)\)。因此 \(q(z)=1/(h(z)+h(a))\) 全纯单叶且模不超过 \(1/r\)。将 \(q\) 乘一个足够小的正数，使像在开单位盘内；接着用圆盘自同构把 \(a\) 的像送到0，再用旋转使导数正实。所得映射属于 \(\mathcal F\)。
\end{proof}
不能直接取 \(1/(z-b)\) 并声称它有界，因为 \(b\) 可能是边界点，\(z\) 可任意接近它。平方根像与负像的分离提供了真正的正距离，而不只是分母处处非零。

在构造有界单叶函数时，取倒数以前必须留下一个真正被像集避开的圆盘。仅知道某点不在像中，不能保证倒数有界；例如 \(1/z\) 在去心平面无零，倒数仍可无界。平方根的像与负像都是开集，又互不相交。在负像中选定一点及其小邻域，便取得整个小圆盘与正像不交的定量间隔。之后才有 \(|1/(h-c)|\le1/\varepsilon\)。这一开集间隔是 Riemann 构造中最容易被“选一个外点”一句话遮住的地方。

\ProseParagraphBreak
还需区分“函数有界”与“像已经是圆盘”。第一次倒数构造只把区域单射地映入一个有界集合。仿射缩放使它成为圆盘内的单叶函数，最后的极值步骤才把像中的遗漏点消去。这样证明并没有事先假设存在满射圆盘坐标，也没有用边界光滑性制造一个预先的映射。

\subsection{正规族紧性}
\label{b3:subsec:050202}

\begin{lemma}\label{b3:lem:05-extremal-compactness}
数 \(S=\sup_{f\in\mathcal F}f'(a)\) 满足 \(0<S<\infty\)。若 \(f_n\in\mathcal F\) 且 \(f_n'(a)\rightarrow S\)，则有局部一致子列，其极限仍属于 \(\mathcal F\)。
\end{lemma}
\begin{proof}
非空性给 \(S>0\)。选 \(\overline{D(a,r)}\subset U\)，全部成员模小于1，所以 Cauchy 估计给 \(f'(a)\le1/r\)，故 \(S\) 有限。Montel 定理提供局部一致子列极限 \(f\)，且 \(|f|\le1\)、\(f(a)=0\)。导数收敛给 \(f'(a)=S>0\)，故它非常数。最大模原理使 \(|f|<1\)，单叶极限定理使 \(f\) 单叶，导数为正实数，因此 \(f\in\mathcal F\)。
\end{proof}
像的边界条件在极限中最容易丢失。单纯的 \(|f_n|<1\) 只给 \(|f|\le1\)；先排除常数退化，才能用最大模把闭盘目标重新缩成开盘。

极值列的导数并非自动有极限。先由 Cauchy 估计给出其有限上界，选取导数趋于上确界的序列，再用 Montel 提取紧集一致收敛子列；导数收敛来自第4章。归一化点处的极限导数为正，使极限不可能常值。单叶函数的紧极限只有“仍单叶”与“常值”两个选择，故正导数恰好排除了丢失单射性的退化。若只要求 \(f(a)=0\) 而不控制导数，\(f_n=f/n\) 就是趋于常值0的合法单叶列。

\ProseParagraphBreak
把函数类中的导数要求为正实数也不是额外限制。任意非零导数都能靠乘一个单位复数转正，而导数的模保持不变。因此最大化正导数与最大化导数模是同一极值问题。这个归一化的好处是極限导数无需在复平面中追踪变化的方向，唯一性证明也不再留有一个未指定旋转。

极值论证对函数族作紧化，而没有对区域边界作紧化。局部有界性使函数在紧集上不会无界，Cauchy 估计进一步控制它们的振荡；对角抽取才把每个局部窗口的极限拼成整域函数。若删去局部有界性，序列 \(nz\) 无法取得有限全纯极限；若保留有界性却删去正导数归一化，极限可能塌成常值。证明的各条件分别阻止不同的退化，因此不能把 Montel 一词视为一个自动生成所需极值函数的黑盒。

\ProseParagraphBreak
导数的上确界也不是由每个函数在区域全体的最大值给出，因为函数模的上确界可为1且不在内部取到。这里最大化的是一个内点的导数，Cauchy 公式把它表示为固定小圆周上的连续线性泛函，紧集一致收敛便保持其值。这个局部量容易在极限中保存，最终却控制了像是否满圆盘，是极值证明尤其有效的原因。

\subsection{极值函数}
\label{b3:subsec:050203}

前一引理给出达到 \(S\) 的函数 \(f\)。极值对象不是区域的面积、边界长度或全域最大导数，而是固定点 \(a\) 处归一化导数的最大值。固定点与相位归一化使比较数是正实数，避免在复数上说“最大导数”。

\ProseParagraphBreak

极值问题的可达性依赖三个封闭步骤：归一化条件由点值与导数收敛保持，单射条件由 Hurwitz 保持，开圆盘目标由非退化与最大模保持。缺少其中任何一项，正规族只给某个极限，并不保证它仍是候选映射。这个论证也解释为什么选择原点导数作目标，而不是一个可能在边界逃逸的全域上确界。

遗漏点的平方根步骤需要谨慎计算它对归一化导数的作用。设当前圆盘像遗漏 \(b\ne0\)，先令 \(g=\phi_b\circ f\)，然后在源区域上取根 \(h^2=g\)。这时 \(|h|<1\)、\(|h(a)|=\sqrt{|b|}\)，但 \(h(a)\) 不再为0。将它送回0所需的自同构，导数模为 \(1/(1-|h(a)|^2)\)。另一方面，平方根链式法则给 \(|h'(a)|=|g'(a)|/(2|h(a)|)\)。合并而非遗漏其中任一因子，才得到改善比
\[
 \frac{1+|b|}{2\sqrt{|b|}}>1.
\]
当遗漏点很靠近0时，这个比值很大；靠近圆周时改善趋近于1，但仍严格大于1。证明所需的只是严格性，不需统一改善量。

\ProseParagraphBreak
这里根函数保持单射，因为根的相同值平方后给 \(g\) 的相同值，\(g\) 又单射。平方根的分支选哪一个不影响导数模，最后的相位归一化还会消除这两个选择。极值论证因此既不要求遗漏集合有面积，也不要求能在遗漏集合中找到一个圆盘；一个遗漏点已足以产生矛盾。和最开始构造有界函数时需要避开的圆盘相比，两步使用外点的方式完全不同。

\subsection{满射性论证}
\label{b3:subsec:050204}

\begin{lemma}[极值映射的满射性]\label{b3:lem:05-extremal-surjective}
达到最大归一化导数的 \(f\in\mathcal F\) 满足 \(f(U)=\mathbb D\)。
\end{lemma}
\begin{proof}
若漏掉 \(b\in\mathbb D\)，因 \(f(a)=0\) 得 \(b\ne0\)。令 \(F=\phi_b\circ f\)，则 \(F\) 无零。区域单连通，存在全纯根 \(h^2=F\)，且 \(|h|<1\)。它单叶，因为 \(h(z_1)=h(z_2)\) 推出 \(F(z_1)=F(z_2)\)，再推出两点相同。令 \(c=h(a)\)，则 \(|c|=\sqrt{|b|}\)。函数 \(g=\phi_c\circ h\) 单叶、取值于圆盘且 \(g(a)=0\)。其导数的模为
\[
 |g'(a)|=\frac{|h'(a)|}{1-|c|^2}
 =\frac{1-|b|^2}{2\sqrt{|b|}(1-|b|)}\cdot f'(a)
 =\frac{1+|b|}{2\sqrt{|b|}}\cdot f'(a).
\]
中间等式使用 \(2h(a)h'(a)=F'(a)\) 与 \(|\phi_b'(0)|=1-|b|^2\)。由于 \(0<|b|<1\)，有 \(1+|b|>2\sqrt{|b|}\)。旋转 \(g\) 使导数正实，就得到 \(\mathcal F\) 中导数大于 \(S\) 的成员，矛盾。
\end{proof}
平方根把遗漏值变成一个可利用的放大因子。它并不将原映射像直接补齐，而是构造一个新的候选映射；极值矛盾保证原映射其实无遗漏。若 \(|b|=1\)，因子等于1，不会产生矛盾，但边界点本来就不是满射到开盘所要求的点。

% !TeX root = ../../Book3.tex
\section{Riemann 映照定理}
\label{b3:sec:0503}

前节的极值构造给出存在性。唯一性则来自两个映射之间的圆盘自同构：把一个映射接在另一个的逆后面，就把区域的形状全部消去，只留下归一化能否固定圆盘坐标的问题。

\subsection{存在性}
\label{b3:subsec:050301}

\term[riemann-yingzhao-dingli]{Riemann 映照定理}[Riemann mapping theorem]在本卷采用真平面区域版本。
\begin{theorem}[Riemann 映照]\label{b3:thm:05-riemann}
设 \(U\subsetneq\mathbb C\) 是单连通区域，\(a\in U\)。则存在唯一双全纯映射 \(f:U\rightarrow\mathbb D\)，满足 \(f(a)=0\)、\(f'(a)>0\)。
\end{theorem}
\begin{proof}
存在性由归一化族非空、极值可达与满射性三引理得到。单叶函数导数非零，所以局部逆全纯；这些局部逆在重叠处因单射一致，拼成全域全纯逆。唯一性见紧随其后的论证。
\end{proof}
存在定理不要求边界光滑，因此即使边界极粗糙，内部共形坐标仍存在；它不自动延拓为欧氏闭包同胚。边界延拓将在\chapref{b3:ch:07}另列条件。

\begin{proposition}\label{b3:prop:05-exhaustion-maps}
设 \(U_1\subset U_2\subset\cdots\) 为真单连通区域，\(a\in U_1\)，且 \(U=\bigcup_n U_n\subsetneq\mathbb C\)。令 \(f_n:U_n\rightarrow\mathbb D\) 和 \(f:U\rightarrow\mathbb D\) 在 \(a\) 按本章归一化。则 \(f_n\) 在 \(U\) 的每个紧集上最终有定义并一致趋于 \(f\)，且 \(r_{U_n}(a)\uparrow r_U(a)\)。
\end{proposition}
\begin{proof}
每个紧集包含于某个 \(U_N\)，于是局部有界与对角抽取给 \(f_n\) 的一个全纯极限 \(g\)。导数单调且不小于 \(f'(a)>0\)，故极限不常值，由单叶极限定理仍单叶并映入圆盘。令 \(h_n=f_n^{-1}\)。函数 \(k_n=f\circ h_n\) 是圆盘自映射且固定0，所以 \(|k_n(z)|\le|z|\)。因此对固定紧盘 \(|z|\le r\)，\(h_n(z)\) 落在 \(f^{-1}(\overline{D(0,r)})\) 内，后者是 \(U\) 内的紧集。对 \(k_n\) 用 Montel，再复合 \(f^{-1}\)，得到 \(h_n\) 的紧极限 \(h:\mathbb D\rightarrow U\)。其导数在0不为零，因为 \(f_n'(a)\) 有正的有限上界。沿同一个子列取极限，在固定紧集上由 \(f_n\circ h_n=\mathrm{id}\) 和 \(h_n\circ f_n=\mathrm{id}\) 得 \(g\circ h=\mathrm{id}\)、\(h\circ g=\mathrm{id}\)。所以 \(g\) 是归一化双全纯映射，唯一性给 \(g=f\)。任何子列都有同一个极限，故全列收敛，导数与共形半径也随之收敛。
\end{proof}
这个命题在逐步逼近复杂区域时很有用，但它只断言内部紧集上的收敛。区域边界可能在逼近过程中改变，映射的边界值不能由此直接交换极限。连通性和单连通性对每个 \(U_n\) 都有要求；任意截取 \(U\cap D(0,n)\) 可能不连通，因此不能未经检查便作为符合假设的穷竭。

\subsection{唯一性}
\label{b3:subsec:050302}

\begin{proof}[\cref{b3:thm:05-riemann}的唯一性证明]
若 \(f,g\) 都满足条件，\(h=g\circ f^{-1}\) 为圆盘自同构且 \(h(0)=0\)。对 \(h\) 和 \(h^{-1}\) 分别用 Schwarz 引理，得到 \(|h'(0)|\le1\) 与 \(|(h^{-1})'(0)|\le1\)。两导数相乘为1，故模都等于1，Schwarz 等号情形给 \(h(z)=cz\)、\(|c|=1\)。另一方面 \(h'(0)=g'(a)/f'(a)>0\)，所以 \(c=1\)，得到 \(g=f\)。
\end{proof}
若只固定 \(f(a)=0\)，仍有一个单位复数旋转自由度；若连 \(a\) 的像都不固定，则还允许所有圆盘自同构。归一化的每一条件分别消除不同自由度，不能把“唯一”理解为所有共形映射完全相同。

\begin{definition}\label{b3:def:05-conformal-radius}
设 \(U\subsetneq\mathbb C\) 单连通，\(a\in U\)，\(F:\mathbb D\rightarrow U\) 双全纯且 \(F(0)=a\)。定义 \(U\) 在 \(a\) 的\term[gongxing-ban-jing]{共形半径}[conformal radius]为 \(r_U(a)=|F'(0)|\)。
\end{definition}
两个这样的逆映射相差圆盘旋转，所以定义与选择无关。若 \(f\) 是本章归一化映射，则 \(r_U(a)=1/f'(a)\)。对圆盘 \(D(c,R)\)，将 \((z-c)/R\) 与圆盘自同构复合得到
\[
 r_{D(c,R)}(a)=\frac{R^2-|a-c|^2}{R}.
\]
对右半平面，归一化坐标可取 \((z-a)/(z+\overline a)\) 再校正相位，故 \(r_U(a)=2\operatorname{Re}a\)。它是一个随点变化的内在尺度，并非到边界的欧氏距离。

\begin{proposition}\label{b3:prop:05-radius-monotonicity}
若 \(U\subset V\subsetneq\mathbb C\) 为单连通区域且 \(a\in U\)，则 \(r_U(a)\le r_V(a)\)，真包含时严格。双全纯 \(T:U\rightarrow W\) 满足 \(r_W(Ta)=|T'(a)|r_U(a)\)。
\end{proposition}
\begin{proof}
取两个归一化映射 \(f_U,f_V\)，则 \(f_V\circ f_U^{-1}\) 把圆盘映入圆盘并固定0。Schwarz 导数估计给 \(f_V'(a)/f_U'(a)\le1\)；取等时该复合是旋转，从而 \(f_V(U)=\mathbb D=f_V(V)\)，单射性迫使 \(U=V\)。共形变换公式由逆映射的链式法则直接得到。
\end{proof}
若 \(D(a,\rho)\subset U\subset D(a,R)\)，便有 \(\rho\le r_U(a)\le R\)。这个简单比较可在不求显式映射的情况下估计归一化导数。后面的 Poincaré 密度将等于 \(2/r_U(a)\)；本章先给出的极值解释会使度量的单调性自然出现。

共形半径还能度量区域对内点的影响。半径 \(R\) 的圆盘中，中心点共形半径为 \(R\)，靠近边界点时却趋于0；所以它描述的是“以当前点为中心的共形尺度”，并非一个区域全局共有的半径。对右半平面，向右平移一点使共形半径线性增大；沿竖直方向移动则保持不变，因为竖直平移是半平面的自同构。这些变化可由显式式子检验，也可由共形协变公式解释。

\ProseParagraphBreak
若区域增加了一条很远的狭窄通道，固定点的共形半径仍可能增加，但仅由集合包含关系只能断言严格单调，不能从图形面积的增加量读出改变量。共形几何的极值量关注全纯函数可以如何变化，而非欧氏面积。在同样面积的两个区域中，共形半径可以不同；在一个无界区域中它也可有限。这正是选取圆盘映射归一化导数作为不变量的意义。

\subsection{圆盘自同构}
\label{b3:subsec:050303}

\begin{proposition}[圆盘自同构]\label{b3:prop:05-disk-automorphisms}
每个圆盘自同构都唯一写成
\[
 h(z)=e^{\mathrm{i}\theta}\frac{z-a}{1-\overline a z},\qquad a\in\mathbb D,
\]
其中 \(e^{\mathrm{i}\theta}\) 是单位复数。反之，这些映射都是圆盘自同构。
\end{proposition}
\begin{proof}
取 \(a=h^{-1}(0)\)，则 \(h\circ\phi_a^{-1}\) 是固定0的圆盘自同构，由刚证明的 Schwarz 等号论证为旋转。反向由第一章的圆盘 Möbius 计算。参数 \(a\) 是唯一零点，旋转因子再由任一点或导数确定。
\end{proof}
圆盘自同构有三个实参数，而 \(f(a)=0\) 给两个实条件，导数正实给一个相位条件。这个参数解释与唯一性证明相符，但参数计数本身不是存在性或唯一性的证明。

\begin{theorem}\label{b3:thm:05-entire-univalent}
整个平面上的单叶整函数恰为 \(f(z)=az+b\)，其中 \(a\ne0\)。
\end{theorem}
\begin{proof}
若无穷远是本性奇点，Casorati--Weierstrass 定理给 \(|z_n|\rightarrow\infty\) 且 \(f(z_n)\rightarrow f(z_0)\)。在 \(z_0\) 附近单叶函数有全纯局部逆，因此充分晚的这些值同时有一个靠近 \(z_0\) 的原像。它与 \(z_n\) 不同，违反单射性。若无穷远可去，Liouville 定理使 \(f\) 常值，亦不可能。故无穷远为极点，\(f\) 是多项式。次数若至少为2，\(f'\) 由代数基本定理有零点，违反单叶函数导数处处非零。反向显然。
\end{proof}
于是全平面也不能双全纯映到任何真平面区域，即使目标无界。这比对有界区域使用 Liouville 更强。证明中的决定性限制是单射性；非常值整函数 \(e^z\) 的像是无界真区域 \(\mathbb C\setminus\{0\}\)，而它确实多对一。

共形类型的唯一性须把双全纯与覆盖区别开。圆盘可通过某些覆盖映到有孔平面域，但这并不说明两域双全纯。双全纯逆映射把任意路径收缩带回源域，因此保存单连通性；圆盘与去心圆盘不能双全纯。指数从平面覆盖去心平面也是多对一映射，与本章平面单叶函数分类不矛盾。第11章引入覆盖后，这种区分将成为统一化陈述的基本组成部分。

\subsection{共形分类}
\label{b3:subsec:050304}

真单连通平面区域因此只有一种共形类型，即圆盘类型。若把全平面也包括进来，平面与圆盘是不同类型，因为 Liouville 排除它们之间的双全纯映射。扩充球面又是第三种，它紧而平面与圆盘不紧；统一化理论将对一般单连通 Riemann 曲面说明这三种模型的穷尽性。

\ProseParagraphBreak

可计算的例子仍然值得保留。右半平面到圆盘可取 \((z-1)/(z+1)\)；水平带 \(|\operatorname{Im}z|<\pi/2\) 先由 \(e^z\) 单叶映到右半平面，再接前式；角域 \(0<\arg z<\alpha\) 先取明确分支 \(z^{\pi/\alpha}\) 映到上半平面，再用 Cayley 变换映到圆盘。每步都需要核对单射与满射，不能仅由边界线的像猜出整个区域的像。

\ProseParagraphBreak

多连通区域有真正的共形不变量。圆环不能映成圆盘，因为双全纯映射是同胚，保持闭曲线是否可收缩。\chapref{b3:ch:08}将以环域模量区分不同圆环；Riemann 映照定理的单连通条件正是在这一位置限制了分类范围。

裂平面的显式映射可以检验抽象证明中的平方根。对 \(U=\mathbb C\setminus(-\infty,0]\)，取 \(\operatorname{Re}\sqrt z>0\) 的平方根分支，再用右半平面的 Cayley 变换，得到
\[
 f(z)=\frac{\sqrt z-1}{\sqrt z+1}.
\]
它在1取0，导数为 \(1/4\)，所以 \(r_U(1)=4\)。点1到裂缝的欧氏距离为1，两种尺度相差显著。裂缝两侧会映到不同的圆周弧，即使它们在平面边界中对应同一负实点；第7章必须以素端处理这种边界辨识。

\ProseParagraphBreak
扇形 \(0<\arg z<\alpha\)、\(0<\alpha\le2\pi\)，可先用 \(z^{\pi/\alpha}=\exp((\pi/\alpha)\log z)\) 映到上半平面，再用分式变换映到圆盘。给定一个明确的辐角区间是必要数据。若忽略分支而直接把幂函数当成全平面函数，所得公式甚至可能不单值。对有理角度，代数式可用于计算；对任意角度，指数定义仍然统一有效。

Riemann 定理给存在性和唯一性，不给一般区域的闭式公式。对于圆盘、半平面、带和扇形，公式来自已知初等函数；对于不规则边界区域，极值构造仍然成立，却可能需要数值逼近才能观察映射。在判断一个候选公式时应逐项核对：在源域单值全纯、导数非零、整体单射、像确为整个目标，以及指定点处的归一化。只画出几个点的像或检查导数，不能替代满射与单射的证明。

\ProseParagraphBreak
例如 \(z\mapsto z^2\) 在右半平面单叶，像为去掉非正实轴的裂平面；若把源域换成上半平面，像则去掉非负实轴。公式相同而像域不同，是因为辐角区间改变。再把源域放宽到全平面，映射变成二对一并在0分歧。所有这些现象都可由辐角范围和局部零点阶数直接核对，体现了“一个函数公式”与“一张区域之间的共形等价”之间的差别。

所有这些显式映射都应在写出逆函数后再核对满射。带的指数逆取对数分支，扇形的幂逆也取指定辐角范围的根。逆的分支范围与正向映射的辐角范围恰好匹配，才保证每个目标点有且只有一个源点。


\begin{chaptersummary}
Riemann 映照定理的存在证明把区域的拓扑条件转化成函数族的极值问题。单连通性保证所需的对数与平方根单值，真区域条件提供一个有限外点，平方根像与其负像的分离提供取倒数所需的共同距离下界。得到有界单叶候选之后，Cauchy 估计控制基点导数，Montel 定理给收敛子列，正的极限导数排除常数退化。最后，任何圆盘内遗漏值都会通过平方根变换产生更大的归一化导数，从而与极值性矛盾。

\ProseParagraphBreak

唯一性不再需要重复存在构造：把两个圆盘坐标复合成自同构，Schwarz 等号情形只留下旋转，导数正实条件再消去这个自由度。共形半径记录逆映射在原点的伸缩，因而同时具有导数极值解释、包含单调性和共形变换公式。对显式映射，必须核对每一步的整体像及逆分支，才能据导数计算这个尺度。对区域穷尽，则必须把观察点放在同一个内部紧集；当极限区域成为全平面时，共形半径可趋于无穷，圆盘坐标随之退化为常数。下一章将继续研究圆盘自映射的导数收缩及其几何含义。
\end{chaptersummary}

\BookExercises

% 本章8题均为自拟；保留原方向，增加基点相位、整体像、周期、等号及退化训练。
\Needspace{6\baselineskip}
\begin{exercise}\label{b3:ex:05-1}
令
\[
 U=\left\{\,w\in\mathbb C\ \middle|\ 
 \operatorname{Re}w>1-\frac{(\operatorname{Im}w)^2}{4}\,\right\}.
\]
用平方根构造 \(U\) 到单位圆盘在 \(w=4\) 处的归一化双全纯映射，并求共形半径。进一步对任意 \(a\in U\)，写出满足 \(f_a(a)=0\)、\(f_a'(a)>0\) 的映射及 \(r_U(a)\)，明确给出所需的相位因子。
\end{exercise}
\begin{solution}
若 \(w\) 为非正实数，则题设不等式要求 \(w>1\)，不可能。因此 \(U\) 不交非正实轴，可以取实部为正的主平方根
\[
 h(w)=\sqrt w=x+\mathrm{i}y,\quad x>0.
\]
写 \(w=u+\mathrm{i}v\)，则 \(u=x^2-y^2\)、\(v=2xy\)，所以
\[
 u-1+\frac{v^2}{4}
 =(x^2-1)(1+y^2).
\]
由于 \(x>0\)，题设不等式恰好等价于 \(x>1\)。这证明 \(h(U)\) 包含于半平面 \(H_1=\{z:\operatorname{Re}z>1\}\)，还须验证反向。任取 \(z=x+\mathrm{i}y\in H_1\)，上式说明 \(z^2\in U\)，且主平方根满足 \(h(z^2)=z\)。所以平方与该分支互为逆映射，\(h:U\rightarrow H_1\) 双全纯。

点4被送到2。将 \(H_1\) 平移为右半平面后使用 Cayley 变换，得到
\[
 f_4(w)=\frac{\sqrt w-2}{\sqrt w}
       =1-\frac2{\sqrt w}.
\]
它的逆映射是
\[
 w=\frac4{(1-\zeta)^2},\quad \zeta\in\mathbb D.
\]
事实上对应的平方根 \(2/(1-\zeta)\) 满足
\[
 \operatorname{Re}\left(\frac2{1-\zeta}-1\right)
 =\frac{1-|\zeta|^2}{|1-\zeta|^2}>0,
\]
故逆像确实落在 \(U\)。又有 \(f_4(4)=0\) 和 \(f_4'(4)=4^{-3/2}=1/8>0\)，所以此式已经完成相位归一化，且
\[
 r_U(4)=8.
\]

对一般基点，记 \(b=h(a)\)、\(x=\operatorname{Re}b>1\)。半平面坐标
\[
 q_a(w)=\frac{h(w)-b}{h(w)+\overline b-2}
\]
把 \(a\) 送到0，并双全纯映到圆盘。其导数为
\[
 q_a'(a)=\frac{h'(a)}{2(x-1)}
        =\frac1{4b(x-1)}.
\]
它通常不是正实数。乘以单位复数 \(b/|b|\) 消去相位，便得到
\[
 f_a(w)=\frac b{|b|}\cdot
 \frac{\sqrt w-b}{\sqrt w+\overline b-2},
 \quad
 f_a'(a)=\frac1{4|b|(x-1)}>0.
\]
因此
\[
 r_U(a)=4|\sqrt a|\left(\operatorname{Re}\sqrt a-1\right).
\]
这里平方根、基点和相位因子必须采用同一个分支；最后代入 \(a=4\) 恢复前面的数值8。
\end{solution}

\Needspace{6\baselineskip}
\begin{exercise}\label{b3:ex:05-2}
把上半单位圆盘
\[
 U=\{\,z\in\mathbb C\mid |z|<1,\ \operatorname{Im}z>0\,\}
\]
双全纯映成单位圆盘，说明每一步的整体像及逆映射所用的辐角范围。进一步计算 \(r_U(\mathrm{i}t)\)，其中 \(0<t<1\)，并检验 \(t\) 接近0或1时的尺度。
\end{exercise}
\begin{solution}
先令
\[
 T(z)=\frac{1+z}{1-z}.
\]
对 \(z\in U\)，直接计算得到
\[
 \operatorname{Re}T(z)=\frac{1-|z|^2}{|1-z|^2}>0,
 \quad
 \operatorname{Im}T(z)=\frac{2\operatorname{Im}z}{|1-z|^2}>0.
\]
因此像包含在第一象限 \(Q\) 中。反之，若 \(w\in Q\)，令 \(z=(w-1)/(w+1)\)。由于
\[
 |w+1|^2-|w-1|^2=4\operatorname{Re}w>0,
 \quad
 \operatorname{Im}z=\frac{2\operatorname{Im}w}{|w+1|^2}>0,
\]
有 \(|z|<1\)、\(\operatorname{Im}z>0\)。所以 \(T\) 双全纯映 \(U\) 到整个 \(Q\)。

平方映射将 \(0<\arg w<\pi/2\) 双射到 \(0<\arg(w^2)<\pi\)，因此把 \(Q\) 双全纯映到上半平面 \(\mathbb H\)。给定 \(\xi=\rho e^{\mathrm{i}\theta}\in\mathbb H\)，唯一的逆像为 \(\sqrt{\rho}\,e^{\mathrm{i}\theta/2}\in Q\)。最后，\((\xi-\mathrm{i})/(\xi+\mathrm{i})\) 把 \(\mathbb H\) 双全纯映到圆盘，逆映射为 \(\mathrm{i}(1+\zeta)/(1-\zeta)\)。于是可取
\[
 F(z)=\frac{T(z)^2-\mathrm{i}}{T(z)^2+\mathrm{i}}.
\]
三步逆映射依次为上半平面 Cayley 逆、辐角落在 \((0,\pi/2)\) 的平方根、以及 \(T^{-1}\)，所以复合的单射与满射都已核对。

计算共形半径时，只需用 \(W=T^2:U\rightarrow\mathbb H\)，无需再调整最终圆盘映射的相位。对 \(a=\mathrm{i}t\)，
\[
 T(\mathrm{i}t)=\frac{1-t^2+2\mathrm{i}t}{1+t^2},
 \quad
 \operatorname{Im}W(\mathrm{i}t)
 =\frac{4t(1-t^2)}{(1+t^2)^2}.
\]
又因 \(|T(\mathrm{i}t)|=1\) 和 \(W'(z)=4T(z)/(1-z)^2\)，有
\[
 |W'(\mathrm{i}t)|=\frac4{1+t^2}.
\]
上半平面在 \(b\) 处的共形半径为 \(2\operatorname{Im}b\)。由共形变换公式
\(r_{\mathbb H}(W(a))=|W'(a)|r_U(a)\)，得到
\[
 r_U(\mathrm{i}t)
 =\frac{2\operatorname{Im}W(\mathrm{i}t)}{|W'(\mathrm{i}t)|}
 =\frac{2t(1-t^2)}{1+t^2}.
\]
例如 \(r_U(\mathrm{i}/2)=3/5\)。当 \(t\downarrow0\) 时，
\(r_U(\mathrm{i}t)/(2t)\to1\)；当 \(t\uparrow1\) 时，
\(r_U(\mathrm{i}t)/(2(1-t))\to1\)。沿虚轴接近直径或上方圆弧时，所得共形尺度分别渐近于到该段边界距离的两倍。
\end{solution}

\Needspace{6\baselineskip}
\begin{exercise}\label{b3:ex:05-3}
在 \(U=\mathbb C\setminus[-1,1]\) 上构造 \(z^2-1\) 的单值全纯平方根 \(h\)，要求在大正实数处取正值，并说明这个条件如何确定分支。证明 \(U\) 非单连通，而且 \(z^2-1\) 与 \(h\) 在 \(U\) 上都没有全纯对数。解释这些结论为何不与平方根的存在冲突。
\end{exercise}
\begin{solution}
考虑
\[
 q(z)=\frac{z-1}{z+1}.
\]
它在 \(U\) 全纯且无零。若 \(z\) 非实，则
\[
 \operatorname{Im}q(z)=\frac{2\operatorname{Im}z}{|z+1|^2}\ne0.
\]
若 \(z\) 为实且不属于 \([-1,1]\)，则 \(q(z)>0\)。所以 \(q(U)\) 避开整个非正实轴，可以使用这一裂平面上的主对数，定义
\[
 h(z)=(z+1)\exp\left(\frac12\operatorname{Log}q(z)\right).
\]
这是一项单值全纯构造，且
\[
 h(z)^2=(z+1)^2q(z)=z^2-1.
\]
对全部实数 \(z>1\)，两个因子均为正实数，故 \(h(z)>0\)。若 \(\widetilde h\) 是另一个全纯根，则 \(\widetilde h/h\) 连续并只取 \(\pm1\)。区域 \(U\) 连通，故该商恒定；在一个大正实数处的正值要求排除负号，确定了唯一分支。

取正向圆周 \(\gamma(t)=Re^{\mathrm{i}t}\)，其中 \(R>1\)。由于 \(1/(z-1)\) 在 \(U\) 全纯，却有
\[
 \int_\gamma\frac{\mathrm{d}z}{z-1}=2\pi\mathrm{i}\ne0,
\]
单连通 Cauchy 定理说明 \(U\) 不可能单连通。再对所考察的两个函数计算对数导数：
\[
 \frac{(z^2-1)'}{z^2-1}
 =\frac1{z-1}+\frac1{z+1},
 \quad
 \frac{h'}h
 =\frac12\left(\frac1{z-1}+\frac1{z+1}\right).
\]
沿 \(\gamma\) 的积分分别为 \(4\pi\mathrm{i}\) 和 \(2\pi\mathrm{i}\)。若相应函数有全纯对数，其对数导数就是该对数的导数，沿闭曲线积分应为零，矛盾。

全纯对数要求所有这样的周期为零；取平方根只要求把周期除以2之后再取指数不产生多值。在这个例子中，环绕线段的一次闭路对 \(z^2-1\) 给出两个零点的总贡献，除以2后的指数仍回到原值。显式的 \(h\) 已经证明平方根存在，而上述非零周期同时证明不能进一步把它写成某个全局对数的指数构造 \(h=e^{L/2}\)、\(e^L=z^2-1\)。某个无零函数有平方根，既不迫使它有对数，也不迫使定义域单连通。
\end{solution}

\Needspace{6\baselineskip}
\begin{exercise}\label{b3:ex:05-4}
设 \(f:\mathbb D\rightarrow\mathbb D\) 与 \(g:\mathbb D\rightarrow\mathbb D\) 全纯，且 \(g\circ f\) 为恒等映射。证明 \(f,g\) 互为圆盘自同构。若只要求左逆 \(g\) 定义于 \(f(\mathbb D)\)，或者允许 \(g:\mathbb D\rightarrow\mathbb C\)，结论是否仍成立？分别说明。
\end{exercise}
\begin{solution}
令 \(a=f(0)\)，则 \(g(a)=0\)。取圆盘自同构
\[
 \phi_a(w)=\frac{w-a}{1-\overline a w},
 \quad
 F=\phi_a\circ f,\quad G=g\circ\phi_a^{-1}.
\]
两者都是圆盘自映射，且 \(F(0)=G(0)=0\)、\(G\circ F=\mathrm{id}_{\mathbb D}\)。由链式法则和 Schwarz 引理，
\[
 G'(0)F'(0)=1,\quad |F'(0)|\le1,\quad |G'(0)|\le1.
\]
因此两个导数模都为1。Schwarz 等号情形给
\(F(z)=\lambda z\)、\(G(z)=\lambda^{-1}z\)，其中 \(|\lambda|=1\)。撤销归一化可知 \(f\) 是圆盘自同构。因为它满射，恒等式 \(g\circ f=\mathrm{id}\) 又在每个圆盘点确定了 \(g=f^{-1}\)，所以 \(g\) 也为自同构。

两个关于左逆的全域条件都有作用。取 \(f(z)=z/2\)，则 \(f\) 单叶却不满射到 \(\mathbb D\)。若左逆只需定义在 \(f(\mathbb D)=D(0,1/2)\)，可取
\[
 g:D(0,1/2)\rightarrow\mathbb D,\quad g(w)=2w.
\]
它确为全纯左逆，但不能作为 \(\mathbb D\) 到自身的映射使用 Schwarz 引理。如果允许目标扩大到 \(\mathbb C\)，同一个公式 \(g(w)=2w\) 可定义在整个 \(\mathbb D\)，仍满足 \(g\circ f=\mathrm{id}\)，但其像不是 \(\mathbb D\) 的子集。因此两种放宽均不能推出 \(f\) 为圆盘自同构。
\end{solution}

\Needspace{6\baselineskip}
\begin{exercise}\label{b3:ex:05-5}
设 \(U,V\subsetneq\mathbb C\) 为单连通区域，\(h:U\rightarrow V\) 全纯。证明对每个 \(a\in U\)，
\[
 |h'(a)|r_U(a)\le r_V(h(a)).
\]
证明在某一点取等，当且仅当 \(h\) 是 \(U\) 到 \(V\) 的双全纯映射；此时在每一点都取等。特别地，取 \(V=\mathbb D\)，推出不要求 \(h(a)=0\) 的估计
\[
 |h'(a)|\le\frac{1-|h(a)|^2}{r_U(a)}.
\]
\end{exercise}
\begin{solution}
固定 \(a\)，记 \(b=h(a)\)。选取分别在 \(a\) 与 \(b\) 处归一化的 Riemann 映射
\[
 f_U:U\rightarrow\mathbb D,\quad
 f_V:V\rightarrow\mathbb D,
\]
使 \(f_U(a)=f_V(b)=0\)，两处导数均为正实数。复合
\[
 H=f_V\circ h\circ f_U^{-1}
\]
是固定原点的圆盘全纯自映射。由归一化和共形半径的定义，
\[
 H'(0)=f_V'(b)h'(a)(f_U^{-1})'(0)
       =\frac{r_U(a)}{r_V(b)}\,h'(a).
\]
Schwarz 引理给 \(|H'(0)|\le1\)，于是得到所求不等式。该论证只要求 \(h\) 全纯，不预先要求单射或满射。

如果在所选点 \(a\) 取等，便有 \(|H'(0)|=1\)。Schwarz 等号情形给 \(H(\zeta)=\lambda\zeta\)、\(|\lambda|=1\)，从而
\[
 h=f_V^{-1}\circ(\lambda f_U)
\]
是双全纯映射。反之，若 \(h\) 双全纯，\(H\) 为固定0的圆盘自同构，故是旋转并达到等号。对任意基点重复这个论证，或者直接使用共形半径的变换公式，便得到处处取等。因而一个点处的等号已经确定了全域可逆性。

单位圆盘的共形半径为 \(r_{\mathbb D}(b)=1-|b|^2\)。代入一般公式即得
\[
 |h'(a)|\le\frac{1-|h(a)|^2}{r_U(a)}.
\]
其等号仍恰好对应于 \(h:U\rightarrow\mathbb D\) 双全纯；基点值没有归零所产生的因子正是 \(1-|h(a)|^2\)，不能把它当成一个固定的常数1。
\end{solution}

\Needspace{6\baselineskip}
\begin{exercise}\label{b3:ex:05-6}
令 \(U_n=D(0,n)\)。在0归一化的圆盘映射为 \(f_n(z)=z/n\)，解释为何它们的极限不能是极限区域的 Riemann 映射。更一般地，设真单连通区域 \(U_n\) 递增穷尽 \(\mathbb C\)，固定 \(a\in U_1\)，令 \(f_n:U_n\rightarrow\mathbb D\) 为在 \(a\) 处归一化的 Riemann 映射。证明
\[
 r_{U_n}(a)\longrightarrow\infty,\quad f_n\longrightarrow0
\]
在全平面的每个固定紧集上按“最终有定义并一致收敛”的意义成立。说明固定阶导数的极限，并判断逆映射 \(f_n^{-1}\) 能否有有限值局部一致收敛子列。
\end{exercise}
\begin{solution}
在最初的圆盘例子中，任意紧集 \(K\subset\mathbb C\) 最终包含于 \(D(0,n)\)，并且
\[
 \sup_{z\in K}|f_n(z)|\le\frac{\max_K|z|}{n}\longrightarrow0.
\]
极限函数是常数，极限区域是全平面。它不满足 Riemann 映照定理中的真区域条件，也不满足正文穷尽命题对极限区域的这一要求。导数 \(1/n\) 的消失正好表现了非退化条件的丢失。

对一般穷尽，固定 \(R>0\)。每个点 \(z\in\overline{D(a,R)}\) 都属于某个 \(U_n\)，这些递增开集覆盖该紧盘。选出有限子覆盖，再取其中指标的最大值，就有
\[
 \overline{D(a,R)}\subset U_N\subset U_n,\quad n\ge N.
\]
圆盘在中心的共形半径为 \(R\)，包含单调性于是给
\(r_{U_n}(a)\ge R\)。因 \(R\) 任意，且这些共形半径随 \(n\) 不减，得到 \(r_{U_n}(a)\to\infty\)，从而 \(f_n'(a)=1/r_{U_n}(a)\to0\)。

仅有基点导数趋于0还没有证明整列函数趋于0。为取得共同紧集上的控制，对 \(n\ge N\) 把 \(f_n\) 限制在 \(D(a,R)\) 上，并对 \(\zeta\mapsto f_n(a+R\zeta)\) 使用 Schwarz 引理，得到
\[
 |f_n(z)|\le\frac{|z-a|}{R},\quad z\in D(a,R).
\]
给定紧集 \(K\subset\mathbb C\) 与误差 \(\varepsilon>0\)，取
\(R>\max_K|z-a|\) 且 \(\max_K|z-a|/R<\varepsilon\)。上述包含关系使全部充分晚的 \(f_n\) 在同一个圆盘上有定义，因此
\[
 \sup_K|f_n|\le\frac{\max_K|z-a|}{R}<\varepsilon.
\]
这证明的是整列在每个固定紧集上趋于0，无须另抽子列。对固定阶导数，先把 \(K\) 扩大为一个紧邻域，再用 Cauchy 估计；扩大后的紧集也最终包含于全部 \(U_n\)，所以每个固定阶导数都在 \(K\) 上一致趋于0。

令 \(G_n=f_n^{-1}:\mathbb D\rightarrow U_n\)。有
\[
 G_n(0)=a,\quad G_n'(0)=r_{U_n}(a)\longrightarrow\infty.
\]
如果某子列在0的一个邻域局部一致趋于有限全纯函数，其原点导数也应收敛到有限值，与上式矛盾。因此逆映射连在0附近都没有有限值局部一致收敛子列。正向坐标压缩为常数和逆向尺度无界，是同一归一化导数关系的两个表现。
\end{solution}

\Needspace{6\baselineskip}
\begin{exercise}\label{b3:ex:05-7}
设 \(m\in\mathbb Z\)，在去心平面 \(\mathbb C^\times\) 上考察无零全纯函数 \(z^m\)。判断它何时有全纯对数；对任意固定整数 \(q\ge2\)，判断它何时有全纯 \(q\) 次根，并列出所有这样的根。证明须通过闭路周期给出必要性，不能只因已经知道一个幂函数候选就断定其他情形不存在。特别说明有平方根与有对数之间的区别。
\end{exercise}
\begin{solution}
若 \(L\) 是 \(z^m\) 的全纯对数，对 \(e^L=z^m\) 求导得 \(L'=m/z\)。沿正向单位圆 \(\gamma\) 积分，
\[
 0=\int_\gamma L'(z)\,\mathrm{d}z
  =m\int_\gamma\frac{\mathrm{d}z}{z}
  =2\pi\mathrm{i}m.
\]
因此必须 \(m=0\)。反之，\(m=0\) 时函数为1，可以取 \(L=0\)；其全部全纯对数为常值 \(2\pi\mathrm{i}k\)，\(k\in\mathbb Z\)，因为连续函数 \(L\) 的值被限制在离散集合 \(2\pi\mathrm{i}\mathbb Z\) 中。

假定全纯 \(q\) 次根 \(g\) 存在。由 \(g^q=z^m\ne0\)，\(g\) 在去心平面无零，且
\[
 q\frac{g'}g=\frac mz.
\]
曲线 \(g\circ\gamma\) 是避开0的闭曲线，其绕0的绕数为某个整数 \(k\)，所以
\[
 \frac1{2\pi\mathrm{i}}\int_\gamma\frac{g'}g\,\mathrm{d}z=k,
 \quad m=qk.
\]
这里整数性来自沿闭曲线连续选取的辐角在终点改变了 \(2\pi\) 的整数倍，不要求 \(g\) 能延拓到0，也不要求 \(g\) 自身有全局对数。因此 \(q\mid m\) 是必要条件，正、负整数 \(m\) 均包括在内。

若 \(m=qk\)，则 \(z^k\) 在 \(\mathbb C^\times\) 全纯，并且 \((z^k)^q=z^m\)。任意另一根与 \(z^k\) 的商只取 \(q\) 次单位根，连通性使该商恒定。所有根恰为
\[
 g(z)=\zeta z^{m/q},\quad \zeta^q=1.
\]
取 \(q=2\)，全纯平方根存在当且仅当 \(m\) 为偶数；取对数却要求 \(m=0\)。例如 \(z^2\) 有根 \(\pm z\)，但没有全纯对数。固定次数的根只要求周期满足相应的整数倍条件，比周期完全消失弱。
\end{solution}

\Needspace{6\baselineskip}
\begin{exercise}\label{b3:ex:05-8}
设 \(F_n:\mathbb C\rightarrow\mathbb C\) 都是单叶整函数。利用正文中的分类结果，证明函数族 \(\{F_n:n\ge1\}\) 在有限值意义下正规，当且仅当两列 \(F_n(0)\)、\(F_n'(0)\) 都有界；刻画其可能的有限子列极限，说明极限何时退化为常数。特别地，若全部成员满足 \(F_n(0)=0\)、\(F_n'(0)=1\)，这个族是什么？

再将“单叶”削弱为“导数处处非零”。构造整函数列 \(g_n\)，仍满足 \(g_n(0)=0\)、\(g_n'(0)=1\)，却在0的任何邻域都没有有限值局部一致收敛子列，并使每个 \(g_n\) 的像是无界真平面区域。核对这一例子为何不违反本章的单叶整函数分类或 Liouville 定理。
\end{exercise}
\begin{solution}
由\cref{b3:thm:05-entire-univalent}，可以且只能写成
\[
 F_n(z)=a_nz+b_n,\quad
 a_n=F_n'(0)\ne0,\quad b_n=F_n(0).
\]
如果 \((a_n)\)、\((b_n)\) 都有界，从族中任取一个序列，先抽取斜率收敛的子列，再在其中抽取截距收敛的子列，得到
\(a_{n_j}\to a\)、\(b_{n_j}\to b\)。对每个固定 \(R>0\)，
\[
 \sup_{|z|\le R}|F_{n_j}(z)-(az+b)|
 \le R|a_{n_j}-a|+|b_{n_j}-b|\longrightarrow0.
\]
因此该族有限值正规。这里允许极限斜率 \(a=0\)，即使每个 \(a_n\ne0\)。

反之，若族有限值正规，基点值或基点导数都不可能有模趋于无穷的子列。否则正规性会从对应函数中再抽出有限值局部一致收敛子列，而点值收敛与导数收敛分别迫使相应参数有有限极限，矛盾。这证明两列参数有界。并且任意有限局部一致子列极限的点值与导数必为 \(b=\lim b_{n_j}\)、\(a=\lim a_{n_j}\)，上述显式式子随即给出极限就是 \(az+b\)。当 \(a\ne0\) 时仍为单叶整函数，当 \(a=0\) 时退化为常数；例如 \(F_n(z)=z/n+b\) 就发生后一种退化。若要求 \(F_n(0)=0\)、\(F_n'(0)=1\)，则每个成员都已等于恒等映射 \(z\)。只把导数规定为正实数并不足以得到正规性，族 \(nz\) 就是反例。

对于削弱后的条件，取
\[
 g_n(z)=\frac{e^{nz}-1}{n}.
\]
直接计算得
\[
 g_n(0)=0,\quad g_n'(0)=1,\quad
 g_n'(z)=e^{nz}\ne0,\quad g_n''(0)=n.
\]
若某子列在0的一个邻域有有限值局部一致极限，第二阶导数在0也必须收敛到有限值，而 \(g_n''(0)=n\) 排除了这种可能。这说明同样的点值和一阶导数归一化，在仅有局部可逆性时不能提供紧性。

最后，指数函数的像为 \(\mathbb C^\times\)，因此
\[
 g_n(\mathbb C)=\mathbb C\setminus\{-1/n\}.
\]
这是无界真区域，不能据“像不等于全平面”套用有界整函数的 Liouville 结论。而 \(g_n(z+2\pi\mathrm{i}/n)=g_n(z)\)，所以各成员都不是单叶函数，也不满足本章分类定理的假设。导数无零、像域为真区域和全局单射之间的区别，在这个同一归一化的函数列中同时显现。
\end{solution}
